\Chapter{120}

\Verse{1}
{
    \zR{话说宝钗听秋纹说袭人不好，连忙进去瞧看，巧姐儿同平儿也随着。}
}
{
    \eR{[English source not present in the checked-in Bencraft/Joly files.]}
}
{
}

\Verse{2}
{
    \zR{走到袭人 炕前，只见袭人心痛难禁，一时气厥。}
}
{
    \eR{[English source not present in the checked-in Bencraft/Joly files.]}
}
{
}

\Verse{3}
{
    \zR{宝钗等用开水灌了过来，仍旧扶他睡下，一 面传请大夫。}
}
{
    \eR{[English source not present in the checked-in Bencraft/Joly files.]}
}
{
}

\Verse{4}
{
    \zR{巧姐儿因问宝钗道：“袭人姐姐怎么病到这个样儿？”}
}
{
    \eR{[English source not present in the checked-in Bencraft/Joly files.]}
}
{
}

\Verse{5}
{
    \zR{宝钗道：“大 前儿晚上哭伤了心了，一时发晕栽倒了。}
}
{
    \eR{[English source not present in the checked-in Bencraft/Joly files.]}
}
{
}

\Verse{6}
{
    \zR{太太叫人扶他回来，他就睡倒了。}
}
{
    \eR{[English source not present in the checked-in Bencraft/Joly files.]}
}
{
}

\Verse{7}
{
    \zR{因外头 有事，没有请大夫瞧他，所以致此。”}
}
{
    \eR{[English source not present in the checked-in Bencraft/Joly files.]}
}
{
}

\Verse{8}
{
    \zR{说着，大夫来了，宝钗等略避。}
}
{
    \eR{[English source not present in the checked-in Bencraft/Joly files.]}
}
{
}

\Verse{9}
{
    \zR{大夫看了脉， 说是急怒所致，开了方子去了。}
}
{
    \eR{[English source not present in the checked-in Bencraft/Joly files.]}
}
{
}

\Verse{10}
{
    \zR{原来袭人模糊听见说宝玉若不回来，便要打发屋里的人都出去，一急越发不好 了。}
}
{
    \eR{[English source not present in the checked-in Bencraft/Joly files.]}
}
{
}

\Verse{11}
{
    \zR{到大夫瞧后，秋纹给他煎药，他各自一人躺着，神魂未定。}
}
{
    \eR{[English source not present in the checked-in Bencraft/Joly files.]}
}
{
}

\Verse{12}
{
    \zR{好像宝玉在他面前， 恍惚又像是见个和尚，手里拿着一本册子揭着看，还说道：“你不是我的人，日后 自然有人家儿的。”}
}
{
    \eR{[English source not present in the checked-in Bencraft/Joly files.]}
}
{
}

\Verse{13}
{
    \zR{袭人似要和他说话，秋纹走来说：“药好了，姐姐吃罢。”}
}
{
    \eR{[English source not present in the checked-in Bencraft/Joly files.]}
}
{
}

\Verse{14}
{
    \zR{袭 人睁眼一瞧，知是个梦，也不告诉人。}
}
{
    \eR{[English source not present in the checked-in Bencraft/Joly files.]}
}
{
}

\Verse{15}
{
    \zR{吃了药，便自己细细的想：“宝玉必是跟了 和尚去。}
}
{
    \eR{[English source not present in the checked-in Bencraft/Joly files.]}
}
{
}

\Verse{16}
{
    \zR{上回他要拿玉出去，便是要脱身的样子。}
}
{
    \eR{[English source not present in the checked-in Bencraft/Joly files.]}
}
{
}

\Verse{17}
{
    \zR{被我揪住，看他竟不像往常，把 我混推混搡的，一点情意都没有。}
}
{
    \eR{[English source not present in the checked-in Bencraft/Joly files.]}
}
{
}

\Verse{18}
{
    \zR{后来待二奶奶更生厌烦，在别的姊妹跟前，也是 没有一点情意：这就是悟道的样子。}
}
{
    \eR{[English source not present in the checked-in Bencraft/Joly files.]}
}
{
}

\Verse{19}
{
    \zR{但是你悟了道，抛了二奶奶怎么好？}
}
{
    \eR{[English source not present in the checked-in Bencraft/Joly files.]}
}
{
}

\Verse{20}
{
    \zR{我是太太 派我服侍你，虽是月钱照着那样的分例，其实我究竟没有在老爷太太跟前回明，就 算了你的屋里人。}
}
{
    \eR{[English source not present in the checked-in Bencraft/Joly files.]}
}
{
}

\Verse{21}
{
    \zR{若是老爷太太打发我出去，我若死守着，又叫人笑话；若是我出 去，心想宝玉待我的情分，实在不忍。”}
}
{
    \eR{[English source not present in the checked-in Bencraft/Joly files.]}
}
{
}

\Verse{22}
{
    \zR{左思右想，万分难处。}
}
{
    \eR{[English source not present in the checked-in Bencraft/Joly files.]}
}
{
}

\Verse{23}
{
    \zR{想到刚才的梦，“说 我是别人的人，那倒不如死了干净。”}
}
{
    \eR{[English source not present in the checked-in Bencraft/Joly files.]}
}
{
}

\Verse{24}
{
    \zR{岂知吃药以后，心痛减了好些，也难躺着， 只好勉强支持。}
}
{
    \eR{[English source not present in the checked-in Bencraft/Joly files.]}
}
{
}

\Verse{25}
{
    \zR{过了几日，起来服侍宝钗。}
}
{
    \eR{[English source not present in the checked-in Bencraft/Joly files.]}
}
{
}

\Verse{26}
{
    \zR{宝钗想念宝玉，暗中垂泪，自叹命苦。}
}
{
    \eR{[English source not present in the checked-in Bencraft/Joly files.]}
}
{
}

\Verse{27}
{
    \zR{又知他母亲打算给哥哥赎罪，很费张罗，不能不帮着打算。}
}
{
    \eR{[English source not present in the checked-in Bencraft/Joly files.]}
}
{
}

\Verse{28}
{
    \zR{暂且不表。}
}
{
    \eR{[English source not present in the checked-in Bencraft/Joly files.]}
}
{
}

\Verse{29}
{
    \zR{且说贾政扶贾母灵柩，贾蓉送了秦氏、凤姐、鸳鸯的棺木到了金陵，先安了葬。}
}
{
    \eR{[English source not present in the checked-in Bencraft/Joly files.]}
}
{
}

\Verse{30}
{
    \zR{贾蓉自送黛玉的灵，也去安葬。}
}
{
    \eR{[English source not present in the checked-in Bencraft/Joly files.]}
}
{
}

\Verse{31}
{
    \zR{贾政料理坟墓的事。}
}
{
    \eR{[English source not present in the checked-in Bencraft/Joly files.]}
}
{
}

\Verse{32}
{
    \zR{一日，接到家书，一行一行的 看到宝玉贾兰得中，心里自是喜欢；后来看到宝玉走失，复又烦恼。}
}
{
    \eR{[English source not present in the checked-in Bencraft/Joly files.]}
}
{
}

\Verse{33}
{
    \zR{只得赶忙回来。}
}
{
    \eR{[English source not present in the checked-in Bencraft/Joly files.]}
}
{
}

\Verse{34}
{
    \zR{在道儿上又闻得有恩赦的旨意，又接着家书，果然赦罪复职，更是喜欢，便日夜趱 行。}
}
{
    \eR{[English source not present in the checked-in Bencraft/Joly files.]}
}
{
}

\Verse{35}
{
    \zR{一日，行到？}
}
{
    \eR{[English source not present in the checked-in Bencraft/Joly files.]}
}
{
}

\Verse{36}
{
    \zR{陵驿地方，那天乍寒，下雪，泊在一个清静去处。}
}
{
    \eR{[English source not present in the checked-in Bencraft/Joly files.]}
}
{
}

\Verse{37}
{
    \zR{贾政打发众人 上岸投帖辞谢朋友，总说即刻开船，都不敢劳动。}
}
{
    \eR{[English source not present in the checked-in Bencraft/Joly files.]}
}
{
}

\Verse{38}
{
    \zR{船上只留一个小厮伺候，自己在 船中写家书，先要打发人起早到家。}
}
{
    \eR{[English source not present in the checked-in Bencraft/Joly files.]}
}
{
}

\Verse{39}
{
    \zR{写到宝玉的事，便停笔。}
}
{
    \eR{[English source not present in the checked-in Bencraft/Joly files.]}
}
{
}

\Verse{40}
{
    \zR{抬头忽见船头上微微 的雪影里面一个人，光着头，赤着脚，身上披着一领大红猩猩毡的斗篷，向贾政倒 身下拜。}
}
{
    \eR{[English source not present in the checked-in Bencraft/Joly files.]}
}
{
}

\Verse{41}
{
    \zR{贾政尚未认清，急忙出船，欲待扶住问他是谁。}
}
{
    \eR{[English source not present in the checked-in Bencraft/Joly files.]}
}
{
}

\Verse{42}
{
    \zR{那人已拜了四拜，站起来 打了个问讯。}
}
{
    \eR{[English source not present in the checked-in Bencraft/Joly files.]}
}
{
}

\Verse{43}
{
    \zR{贾政才要还揖，迎面一看，不是别人，却是宝玉。}
}
{
    \eR{[English source not present in the checked-in Bencraft/Joly files.]}
}
{
}

\Verse{44}
{
    \zR{贾政吃一大惊，忙 问道：“可是宝玉么？”}
}
{
    \eR{[English source not present in the checked-in Bencraft/Joly files.]}
}
{
}

\Verse{45}
{
    \zR{那人只不言语，似喜似悲。}
}
{
    \eR{[English source not present in the checked-in Bencraft/Joly files.]}
}
{
}

\Verse{46}
{
    \zR{贾政又问道：“你若是宝玉， 如何这样打扮，跑到这里来？”}
}
{
    \eR{[English source not present in the checked-in Bencraft/Joly files.]}
}
{
}

\Verse{47}
{
    \zR{宝玉未及回言，只见船头上来了两人，一僧一道， 夹住宝玉道：“俗缘已毕，还不快走。”}
}
{
    \eR{[English source not present in the checked-in Bencraft/Joly files.]}
}
{
}

\Verse{48}
{
    \zR{说着，三个人飘然登岸而去。}
}
{
    \eR{[English source not present in the checked-in Bencraft/Joly files.]}
}
{
}

\Verse{49}
{
    \zR{贾政不顾地 滑，疾忙来赶，见那三人在前，那里赶得上？}
}
{
    \eR{[English source not present in the checked-in Bencraft/Joly files.]}
}
{
}

\Verse{50}
{
    \zR{只听得他们三人口中不知是那个作歌 曰： 我所居兮青埂之峰，我所游兮鸿蒙太空。}
}
{
    \eR{[English source not present in the checked-in Bencraft/Joly files.]}
}
{
}

\Verse{51}
{
    \zR{谁与我逝兮吾谁与从？}
}
{
    \eR{[English source not present in the checked-in Bencraft/Joly files.]}
}
{
}

\Verse{52}
{
    \zR{渺渺茫茫兮归彼大荒! 贾政一面听着，一面赶去，转过一小坡，倏然不见。}
}
{
    \eR{[English source not present in the checked-in Bencraft/Joly files.]}
}
{
}

\Verse{53}
{
    \zR{贾政已赶得心虚气喘，惊疑不 定。}
}
{
    \eR{[English source not present in the checked-in Bencraft/Joly files.]}
}
{
}

\Verse{54}
{
    \zR{回过头来，见自己的小厮也随后赶来，贾政问道：“你看见方才那三个人么？”}
}
{
    \eR{[English source not present in the checked-in Bencraft/Joly files.]}
}
{
}

\Verse{55}
{
    \zR{小厮道：“看见的。}
}
{
    \eR{[English source not present in the checked-in Bencraft/Joly files.]}
}
{
}

\Verse{56}
{
    \zR{奴才为老爷追赶，故也赶来。}
}
{
    \eR{[English source not present in the checked-in Bencraft/Joly files.]}
}
{
}

\Verse{57}
{
    \zR{后来只见老爷，不见那三个人了。”}
}
{
    \eR{[English source not present in the checked-in Bencraft/Joly files.]}
}
{
}

\Verse{58}
{
    \zR{贾政还欲前走，只见白茫茫一片旷野，并无一人。}
}
{
    \eR{[English source not present in the checked-in Bencraft/Joly files.]}
}
{
}

\Verse{59}
{
    \zR{贾政知是古怪，只得回来。}
}
{
    \eR{[English source not present in the checked-in Bencraft/Joly files.]}
}
{
}

\Verse{60}
{
    \zR{众家人回船，见贾政不在舱中，问了船夫，说是老爷上岸追赶两个和尚一个道 士去了。}
}
{
    \eR{[English source not present in the checked-in Bencraft/Joly files.]}
}
{
}

\Verse{61}
{
    \zR{众人也从雪地里寻踪迎去，远远见贾政来了，迎上去接着，一同回船。}
}
{
    \eR{[English source not present in the checked-in Bencraft/Joly files.]}
}
{
}

\Verse{62}
{
    \zR{贾 政坐下，喘息方定，将见宝玉的话说了一遍。}
}
{
    \eR{[English source not present in the checked-in Bencraft/Joly files.]}
}
{
}

\Verse{63}
{
    \zR{众人回禀，便要在这地方寻觅。}
}
{
    \eR{[English source not present in the checked-in Bencraft/Joly files.]}
}
{
}

\Verse{64}
{
    \zR{贾政 叹道：“你们不知道，这是我亲眼见的，并非鬼怪。}
}
{
    \eR{[English source not present in the checked-in Bencraft/Joly files.]}
}
{
}

\Verse{65}
{
    \zR{况听得歌声，大有玄妙。}
}
{
    \eR{[English source not present in the checked-in Bencraft/Joly files.]}
}
{
}

\Verse{66}
{
    \zR{宝玉 生下时，衔了玉来，便也古怪，我早知是不祥之兆，为的是老太太疼爱，所以养育 到今。}
}
{
    \eR{[English source not present in the checked-in Bencraft/Joly files.]}
}
{
}

\Verse{67}
{
    \zR{便是那和尚道士，我也见了三次：头一次，是那僧道来说玉的好处；第二次， 便是宝玉病重，他来了，将那玉持诵了一番，宝玉便好了；第三次，送那玉来，坐
        在前厅，我一转眼就不见了。}
}
{
    \eR{[English source not present in the checked-in Bencraft/Joly files.]}
}
{
}

\Verse{68}
{
    \zR{我心里便有些诧异，只道宝玉果真有造化，高僧仙道 来护佑他的。}
}
{
    \eR{[English source not present in the checked-in Bencraft/Joly files.]}
}
{
}

\Verse{69}
{
    \zR{岂知宝玉是下凡历劫的，竟哄了老太太十九年!如今叫我才明白。”}
}
{
    \eR{[English source not present in the checked-in Bencraft/Joly files.]}
}
{
}

\Verse{70}
{
    \zR{说到那里，掉下泪来。}
}
{
    \eR{[English source not present in the checked-in Bencraft/Joly files.]}
}
{
}

\Verse{71}
{
    \zR{众人道：“宝二爷果然是下凡的和尚，就不该中举人了。}
}
{
    \eR{[English source not present in the checked-in Bencraft/Joly files.]}
}
{
}

\Verse{72}
{
    \zR{怎 么中了才去？”}
}
{
    \eR{[English source not present in the checked-in Bencraft/Joly files.]}
}
{
}

\Verse{73}
{
    \zR{贾政道：“你们那里知道？}
}
{
    \eR{[English source not present in the checked-in Bencraft/Joly files.]}
}
{
}

\Verse{74}
{
    \zR{大凡天上星宿，山中老僧，洞里的精灵， 他自具一种性情。}
}
{
    \eR{[English source not present in the checked-in Bencraft/Joly files.]}
}
{
}

\Verse{75}
{
    \zR{你看宝玉何尝肯念书？}
}
{
    \eR{[English source not present in the checked-in Bencraft/Joly files.]}
}
{
}

\Verse{76}
{
    \zR{他若略一经心，无有不能的。}
}
{
    \eR{[English source not present in the checked-in Bencraft/Joly files.]}
}
{
}

\Verse{77}
{
    \zR{他那一种脾 气，也是各别另样。”}
}
{
    \eR{[English source not present in the checked-in Bencraft/Joly files.]}
}
{
}

\Verse{78}
{
    \zR{说着又叹了几声。}
}
{
    \eR{[English source not present in the checked-in Bencraft/Joly files.]}
}
{
}

\Verse{79}
{
    \zR{众人便拿兰哥得中、家道复兴的话解了一 番。}
}
{
    \eR{[English source not present in the checked-in Bencraft/Joly files.]}
}
{
}

\Verse{80}
{
    \zR{贾政仍旧写家书，便把这事写上，劝谕合家不必想念了。}
}
{
    \eR{[English source not present in the checked-in Bencraft/Joly files.]}
}
{
}

\Verse{81}
{
    \zR{写完封好，即着家人 回去，贾政随后赶回。}
}
{
    \eR{[English source not present in the checked-in Bencraft/Joly files.]}
}
{
}

\Verse{82}
{
    \zR{暂且不提。}
}
{
    \eR{[English source not present in the checked-in Bencraft/Joly files.]}
}
{
}

\Verse{83}
{
    \zR{且说薛姨妈得了赦罪的信，便命薛蝌去各处借贷，并自己凑齐了赎罪银两。}
}
{
    \eR{[English source not present in the checked-in Bencraft/Joly files.]}
}
{
}

\Verse{84}
{
    \zR{刑 部准了，收兑了银子，一角文书，将薛蟠放出。}
}
{
    \eR{[English source not present in the checked-in Bencraft/Joly files.]}
}
{
}

\Verse{85}
{
    \zR{他们母子姊妹弟兄见面，不必细述， 自然是悲喜交集了。}
}
{
    \eR{[English source not present in the checked-in Bencraft/Joly files.]}
}
{
}

\Verse{86}
{
    \zR{薛蟠自己立誓说道：“若是再犯前病，必定犯杀犯剐！”}
}
{
    \eR{[English source not present in the checked-in Bencraft/Joly files.]}
}
{
}

\Verse{87}
{
    \zR{薛姨 妈见他这样，便握他的嘴，说：“只要自己拿定主意，必定还要妄口巴舌血淋淋的 起这样恶誓么？}
}
{
    \eR{[English source not present in the checked-in Bencraft/Joly files.]}
}
{
}

\Verse{88}
{
    \zR{只是香菱跟你受了多少苦处，你媳妇儿已经自己治死自己了，如今 虽说穷了，这碗饭还有得吃：据我的主意，我便算他是媳妇了。}
}
{
    \eR{[English source not present in the checked-in Bencraft/Joly files.]}
}
{
}

\Verse{89}
{
    \zR{你心里怎么样？”}
}
{
    \eR{[English source not present in the checked-in Bencraft/Joly files.]}
}
{
}

\Verse{90}
{
    \zR{薛蟠点头愿意。}
}
{
    \eR{[English source not present in the checked-in Bencraft/Joly files.]}
}
{
}

\Verse{91}
{
    \zR{宝钗等也说：“很该这样。”}
}
{
    \eR{[English source not present in the checked-in Bencraft/Joly files.]}
}
{
}

\Verse{92}
{
    \zR{倒把香菱急得脸胀通红，说是：“伏 侍大爷一样的，何必如此？”}
}
{
    \eR{[English source not present in the checked-in Bencraft/Joly files.]}
}
{
}

\Verse{93}
{
    \zR{众人便称起“大奶奶”来，无人不服。}
}
{
    \eR{[English source not present in the checked-in Bencraft/Joly files.]}
}
{
}

\Verse{94}
{
    \zR{薛蟠便要去拜谢贾家。}
}
{
    \eR{[English source not present in the checked-in Bencraft/Joly files.]}
}
{
}

\Verse{95}
{
    \zR{薛姨妈宝钗也都过来。}
}
{
    \eR{[English source not present in the checked-in Bencraft/Joly files.]}
}
{
}

\Verse{96}
{
    \zR{见了众人，彼此聚首，又说了一 番的话。}
}
{
    \eR{[English source not present in the checked-in Bencraft/Joly files.]}
}
{
}

\Verse{97}
{
    \zR{正说着，恰好那日贾政的家人回家，呈上书子，说：“老爷不日到了。”}
}
{
    \eR{[English source not present in the checked-in Bencraft/Joly files.]}
}
{
}

\Verse{98}
{
    \zR{王夫人叫贾兰将书子念给听。}
}
{
    \eR{[English source not present in the checked-in Bencraft/Joly files.]}
}
{
}

\Verse{99}
{
    \zR{贾兰念到贾政亲见宝玉的一段，众人听了，都痛哭起 来，王夫人、宝钗、袭人等更甚。}
}
{
    \eR{[English source not present in the checked-in Bencraft/Joly files.]}
}
{
}

\Verse{100}
{
    \zR{大家又将贾政书内叫家内不必悲伤，原是借胎的 话解说了一番：“与其作了官，倘或命运不好，犯了事，坏家败产，那时倒不好了。}
}
{
    \eR{[English source not present in the checked-in Bencraft/Joly files.]}
}
{
}

\Verse{101}
{
    \zR{宁可咱们家出一位佛爷，倒是老爷太太的积德，所以才投到咱们家来。}
}
{
    \eR{[English source not present in the checked-in Bencraft/Joly files.]}
}
{
}

\Verse{102}
{
    \zR{不是说句不 顾前后的话：当初东府里太爷，倒是修炼了十几年，也没有成了仙，这佛是更难成 的。}
}
{
    \eR{[English source not present in the checked-in Bencraft/Joly files.]}
}
{
}

\Verse{103}
{
    \zR{太太这么一想，心里便开豁了。”}
}
{
    \eR{[English source not present in the checked-in Bencraft/Joly files.]}
}
{
}

\Verse{104}
{
    \zR{王夫人哭着和薛姨妈道：“宝玉抛了我，我 还恨他呢。}
}
{
    \eR{[English source not present in the checked-in Bencraft/Joly files.]}
}
{
}

\Verse{105}
{
    \zR{我叹的是媳妇的命苦，才成了一二年的亲，怎么他就硬着肠子，都撂下 了走了呢！”}
}
{
    \eR{[English source not present in the checked-in Bencraft/Joly files.]}
}
{
}

\Verse{106}
{
    \zR{薛姨妈听了，也甚伤心。}
}
{
    \eR{[English source not present in the checked-in Bencraft/Joly files.]}
}
{
}

\Verse{107}
{
    \zR{宝钗哭得人事不知。}
}
{
    \eR{[English source not present in the checked-in Bencraft/Joly files.]}
}
{
}

\Verse{108}
{
    \zR{所有爷们都在外头。}
}
{
    \eR{[English source not present in the checked-in Bencraft/Joly files.]}
}
{
}

\Verse{109}
{
    \zR{王夫人便说道：“我为他担了一辈子 的惊，刚刚儿的娶了亲，中了举人，又知道媳妇作了胎，我才喜欢些，不想弄到这
        样结局!早知这样，就不该娶亲，害了人家的姑娘。”}
}
{
    \eR{[English source not present in the checked-in Bencraft/Joly files.]}
}
{
}

\Verse{110}
{
    \zR{薛姨妈道：“这是自己一定 的。}
}
{
    \eR{[English source not present in the checked-in Bencraft/Joly files.]}
}
{
}

\Verse{111}
{
    \zR{咱们这样人家，还有什么别的说的吗？}
}
{
    \eR{[English source not present in the checked-in Bencraft/Joly files.]}
}
{
}

\Verse{112}
{
    \zR{幸喜有了胎，将来生个外孙子，必定是 有成立的，后来就有了结果了。}
}
{
    \eR{[English source not present in the checked-in Bencraft/Joly files.]}
}
{
}

\Verse{113}
{
    \zR{你看大奶奶，如今兰哥儿中了举人，明年成了进士， 可不是就做了官了么？}
}
{
    \eR{[English source not present in the checked-in Bencraft/Joly files.]}
}
{
}

\Verse{114}
{
    \zR{他头里的苦也算吃尽的了，如今的甜来，也是他为人的好处。}
}
{
    \eR{[English source not present in the checked-in Bencraft/Joly files.]}
}
{
}

\Verse{115}
{
    \zR{我们姑娘的心肠儿姐姐是知道的，并不是刻薄轻佻的人，姐姐倒不必耽忧。”}
}
{
    \eR{[English source not present in the checked-in Bencraft/Joly files.]}
}
{
}

\Verse{116}
{
    \zR{王夫 人被薛姨妈一番言语说得极有理，心想：“宝钗小时候便是廉静寡欲极爱素淡的， 他所以才有这个事。}
}
{
    \eR{[English source not present in the checked-in Bencraft/Joly files.]}
}
{
}

\Verse{117}
{
    \zR{想人生在世，真有个定数的。}
}
{
    \eR{[English source not present in the checked-in Bencraft/Joly files.]}
}
{
}

\Verse{118}
{
    \zR{看着宝钗虽是痛哭，他那端庄样 儿一点不走，却倒来劝我，这是真真难得。}
}
{
    \eR{[English source not present in the checked-in Bencraft/Joly files.]}
}
{
}

\Verse{119}
{
    \zR{不想宝玉这样一个人，红尘中福分竟没 有一点儿！”}
}
{
    \eR{[English source not present in the checked-in Bencraft/Joly files.]}
}
{
}

\Verse{120}
{
    \zR{想了一回，也觉解了好些。}
}
{
    \eR{[English source not present in the checked-in Bencraft/Joly files.]}
}
{
}

\Verse{121}
{
    \zR{又想到袭人身上：“若说别的丫头呢，没 有什么难处的：大的配了出去，小的伏侍二奶奶就是了。}
}
{
    \eR{[English source not present in the checked-in Bencraft/Joly files.]}
}
{
}

\Verse{122}
{
    \zR{独有袭人可怎么处呢？”}
}
{
    \eR{[English source not present in the checked-in Bencraft/Joly files.]}
}
{
}

\Verse{123}
{
    \zR{此时人多也不好说，且等晚上和薛姨妈商量。}
}
{
    \eR{[English source not present in the checked-in Bencraft/Joly files.]}
}
{
}

\Verse{124}
{
    \zR{那日薛姨妈并未回家，因恐宝钗痛哭，住在宝钗房中解劝。}
}
{
    \eR{[English source not present in the checked-in Bencraft/Joly files.]}
}
{
}

\Verse{125}
{
    \zR{那宝钗却是极明理， 思前想后：“宝玉原是一种奇异的人，夙世前因，自有一定，原无可怨天尤人。”}
}
{
    \eR{[English source not present in the checked-in Bencraft/Joly files.]}
}
{
}

\Verse{126}
{
    \zR{更将大道理的话告诉他母亲了。}
}
{
    \eR{[English source not present in the checked-in Bencraft/Joly files.]}
}
{
}

\Verse{127}
{
    \zR{薛姨妈心里反倒安慰，便到王夫人那里，先把宝钗 的话说了。}
}
{
    \eR{[English source not present in the checked-in Bencraft/Joly files.]}
}
{
}

\Verse{128}
{
    \zR{王夫人点头叹道：“若说我无德，不该有这样好媳妇了。”}
}
{
    \eR{[English source not present in the checked-in Bencraft/Joly files.]}
}
{
}

\Verse{129}
{
    \zR{说着更又伤 心起来。}
}
{
    \eR{[English source not present in the checked-in Bencraft/Joly files.]}
}
{
}

\Verse{130}
{
    \zR{薛姨妈倒又劝了一会子。}
}
{
    \eR{[English source not present in the checked-in Bencraft/Joly files.]}
}
{
}

\Verse{131}
{
    \zR{因又提起袭人来，说：“我见袭人近来瘦的了不 得，他是一心想着宝哥儿。}
}
{
    \eR{[English source not present in the checked-in Bencraft/Joly files.]}
}
{
}

\Verse{132}
{
    \zR{但是正配呢理应守的，屋里人愿守也是有的。}
}
{
    \eR{[English source not present in the checked-in Bencraft/Joly files.]}
}
{
}

\Verse{133}
{
    \zR{惟有这袭 人，虽说是算个屋里人，到底他和宝哥儿并没有过明路儿的。”}
}
{
    \eR{[English source not present in the checked-in Bencraft/Joly files.]}
}
{
}

\Verse{134}
{
    \zR{王夫人道：“我才 刚想着，正要等妹妹商量商量。}
}
{
    \eR{[English source not present in the checked-in Bencraft/Joly files.]}
}
{
}

\Verse{135}
{
    \zR{若说放他出去，恐怕他不愿意，又要寻死觅活的； 若要留着他也罢，又恐老爷不依：所以难处。”}
}
{
    \eR{[English source not present in the checked-in Bencraft/Joly files.]}
}
{
}

\Verse{136}
{
    \zR{薛姨妈道：“我看姨老爷是再不肯 叫守着的。}
}
{
    \eR{[English source not present in the checked-in Bencraft/Joly files.]}
}
{
}

\Verse{137}
{
    \zR{再者，姨老爷并不知道袭人的事，想来不过是个丫头，那有留的理呢？}
}
{
    \eR{[English source not present in the checked-in Bencraft/Joly files.]}
}
{
}

\Verse{138}
{
    \zR{只要姐姐叫他本家的人来，狠狠的吩咐他，叫他配一门正经亲事，再多多的陪送他 些东西。}
}
{
    \eR{[English source not present in the checked-in Bencraft/Joly files.]}
}
{
}

\Verse{139}
{
    \zR{那孩子心肠儿也好，年纪儿又轻，也不枉跟了姐姐会子，也算姐姐待他不 薄了。}
}
{
    \eR{[English source not present in the checked-in Bencraft/Joly files.]}
}
{
}

\Verse{140}
{
    \zR{袭人那里，还得我细细劝他。}
}
{
    \eR{[English source not present in the checked-in Bencraft/Joly files.]}
}
{
}

\Verse{141}
{
    \zR{就是叫他家的人来，也不用告诉他；只等他家 里果然说定了好人家儿，我们还打听打听，若果然足衣足食、女婿长的像个人儿， 然后叫他出去。”}
}
{
    \eR{[English source not present in the checked-in Bencraft/Joly files.]}
}
{
}

\Verse{142}
{
    \zR{王夫人听了，道：“这个主意很是。}
}
{
    \eR{[English source not present in the checked-in Bencraft/Joly files.]}
}
{
}

\Verse{143}
{
    \zR{不然叫老爷冒冒失失的一办， 我可不是又害了一个人了么？”}
}
{
    \eR{[English source not present in the checked-in Bencraft/Joly files.]}
}
{
}

\Verse{144}
{
    \zR{薛姨妈听了，点头道：“可不是么？”}
}
{
    \eR{[English source not present in the checked-in Bencraft/Joly files.]}
}
{
}

\Verse{145}
{
    \zR{又说了几句， 便辞了王夫人仍到宝钗房中去了。}
}
{
    \eR{[English source not present in the checked-in Bencraft/Joly files.]}
}
{
}

\Verse{146}
{
    \zR{看见袭人泪痕满面，薛姨妈便劝解譬喻了一会。}
}
{
    \eR{[English source not present in the checked-in Bencraft/Joly files.]}
}
{
}

\Verse{147}
{
    \zR{袭人本来老实，不是伶牙俐齿的人，薛姨妈说一句，他应一句，回来说道：“我是 做下人的人，姨太太瞧得起我，才和我说这些话。}
}
{
    \eR{[English source not present in the checked-in Bencraft/Joly files.]}
}
{
}

\Verse{148}
{
    \zR{我是从不敢违拗太太的。”}
}
{
    \eR{[English source not present in the checked-in Bencraft/Joly files.]}
}
{
}

\Verse{149}
{
    \zR{薛姨 妈听他的话，“好一个柔顺的孩子！”}
}
{
    \eR{[English source not present in the checked-in Bencraft/Joly files.]}
}
{
}

\Verse{150}
{
    \zR{心里更加喜欢。}
}
{
    \eR{[English source not present in the checked-in Bencraft/Joly files.]}
}
{
}

\Verse{151}
{
    \zR{宝钗又将大义的话说了一遍， 大家各自相安。}
}
{
    \eR{[English source not present in the checked-in Bencraft/Joly files.]}
}
{
}

\Verse{152}
{
    \zR{过了几日，贾政回家，众人迎接。}
}
{
    \eR{[English source not present in the checked-in Bencraft/Joly files.]}
}
{
}

\Verse{153}
{
    \zR{贾政见贾赦贾珍已都回家，弟兄叔侄相见， 大家历叙别来的景况。}
}
{
    \eR{[English source not present in the checked-in Bencraft/Joly files.]}
}
{
}

\Verse{154}
{
    \zR{然后内眷们见了，不免想起宝玉来，又大家伤了一会子心。}
}
{
    \eR{[English source not present in the checked-in Bencraft/Joly files.]}
}
{
}

\Verse{155}
{
    \zR{贾政喝住道：“这是一定的道理!如今只要我们在外把持家事，你们在内相助，断 不可仍是从前这样的散漫。}
}
{
    \eR{[English source not present in the checked-in Bencraft/Joly files.]}
}
{
}

\Verse{156}
{
    \zR{别房的事，各有各家料理，也不用承总。}
}
{
    \eR{[English source not present in the checked-in Bencraft/Joly files.]}
}
{
}

\Verse{157}
{
    \zR{我们本房的事， 里头全归于你，都要按理而行。”}
}
{
    \eR{[English source not present in the checked-in Bencraft/Joly files.]}
}
{
}

\Verse{158}
{
    \zR{王夫人便将宝钗有孕的话也告诉了，“将来丫头 们都放出去。”}
}
{
    \eR{[English source not present in the checked-in Bencraft/Joly files.]}
}
{
}

\Verse{159}
{
    \zR{贾政听了，点头无语。}
}
{
    \eR{[English source not present in the checked-in Bencraft/Joly files.]}
}
{
}

\Verse{160}
{
    \zR{次日，贾政进内请示大臣们，说是：“蒙恩感激。}
}
{
    \eR{[English source not present in the checked-in Bencraft/Joly files.]}
}
{
}

\Verse{161}
{
    \zR{但未服阕，应该怎么谢恩之 处，望乞大人们指教。”}
}
{
    \eR{[English source not present in the checked-in Bencraft/Joly files.]}
}
{
}

\Verse{162}
{
    \zR{众朝臣说是代奏请旨。}
}
{
    \eR{[English source not present in the checked-in Bencraft/Joly files.]}
}
{
}

\Verse{163}
{
    \zR{于是圣恩浩荡，即命陛见。}
}
{
    \eR{[English source not present in the checked-in Bencraft/Joly files.]}
}
{
}

\Verse{164}
{
    \zR{贾政进 内谢了恩。}
}
{
    \eR{[English source not present in the checked-in Bencraft/Joly files.]}
}
{
}

\Verse{165}
{
    \zR{圣上又降了好些旨意，又问起宝玉的事来。}
}
{
    \eR{[English source not present in the checked-in Bencraft/Joly files.]}
}
{
}

\Verse{166}
{
    \zR{贾政据实回奏。}
}
{
    \eR{[English source not present in the checked-in Bencraft/Joly files.]}
}
{
}

\Verse{167}
{
    \zR{圣上称奇， 旨意说：宝玉的文章固是清奇，想他必是过来人，所以如此。}
}
{
    \eR{[English source not present in the checked-in Bencraft/Joly files.]}
}
{
}

\Verse{168}
{
    \zR{若在朝中，可以进用； 他既不敢受圣朝的爵位，便赏了一个“文妙真人”的道号。}
}
{
    \eR{[English source not present in the checked-in Bencraft/Joly files.]}
}
{
}

\Verse{169}
{
    \zR{贾政又叩头谢恩而出。}
}
{
    \eR{[English source not present in the checked-in Bencraft/Joly files.]}
}
{
}

\Verse{170}
{
    \zR{回到家中，贾琏贾珍接着，贾政将朝内的话述了一遍，众人喜欢。}
}
{
    \eR{[English source not present in the checked-in Bencraft/Joly files.]}
}
{
}

\Verse{171}
{
    \zR{贾珍便回说：“宁 国府第，收拾齐全，回明了要搬过去。}
}
{
    \eR{[English source not present in the checked-in Bencraft/Joly files.]}
}
{
}

\Verse{172}
{
    \zR{栊翠庵圈在园内，给四妹妹养静。”}
}
{
    \eR{[English source not present in the checked-in Bencraft/Joly files.]}
}
{
}

\Verse{173}
{
    \zR{贾政并 不言语，隔了半日，却吩咐了一番仰报天恩的话。}
}
{
    \eR{[English source not present in the checked-in Bencraft/Joly files.]}
}
{
}

\Verse{174}
{
    \zR{贾琏也趁便回说：“巧姐亲事，父亲太太都愿意给周家为媳。”}
}
{
    \eR{[English source not present in the checked-in Bencraft/Joly files.]}
}
{
}

\Verse{175}
{
    \zR{贾政昨晚也知 巧姐的始末，便说：“大老爷大太太作主就是了。}
}
{
    \eR{[English source not present in the checked-in Bencraft/Joly files.]}
}
{
}

\Verse{176}
{
    \zR{莫说村居不好，只要人家清白， 孩子肯念书，能够上进。}
}
{
    \eR{[English source not present in the checked-in Bencraft/Joly files.]}
}
{
}

\Verse{177}
{
    \zR{朝里那些官，难道都是城里的人么？”}
}
{
    \eR{[English source not present in the checked-in Bencraft/Joly files.]}
}
{
}

\Verse{178}
{
    \zR{贾琏答应了“是”， 又说：“父亲有了年纪，况且又有痰症的根子，静养几年，诸事原仗二老爷为主。”}
}
{
    \eR{[English source not present in the checked-in Bencraft/Joly files.]}
}
{
}

\Verse{179}
{
    \zR{贾政道：“提起村居养静，甚合我意，只是我受恩深重，尚未酬报耳。”}
}
{
    \eR{[English source not present in the checked-in Bencraft/Joly files.]}
}
{
}

\Verse{180}
{
    \zR{贾政说毕 进内，贾琏打发请了刘老老来，应了这件事。}
}
{
    \eR{[English source not present in the checked-in Bencraft/Joly files.]}
}
{
}

\Verse{181}
{
    \zR{刘老老见了王夫人等，便说些将来怎 样升官，怎样起家，怎样子孙昌盛。}
}
{
    \eR{[English source not present in the checked-in Bencraft/Joly files.]}
}
{
}

\Verse{182}
{
    \zR{正说着，丫头回道：“花自芳的女人进来请安。”}
}
{
    \eR{[English source not present in the checked-in Bencraft/Joly files.]}
}
{
}

\Verse{183}
{
    \zR{王夫人问几句话，花自芳的 女人将亲戚作媒，说的是城南蒋家的，现在有房有地，又有铺面。}
}
{
    \eR{[English source not present in the checked-in Bencraft/Joly files.]}
}
{
}

\Verse{184}
{
    \zR{姑爷年纪略大几 岁，并没有娶过的，况且人物儿长的是百里挑一的。}
}
{
    \eR{[English source not present in the checked-in Bencraft/Joly files.]}
}
{
}

\Verse{185}
{
    \zR{王夫人听了愿意，说道：“你 去应了，隔几日进来，再接你妹子罢。”}
}
{
    \eR{[English source not present in the checked-in Bencraft/Joly files.]}
}
{
}

\Verse{186}
{
    \zR{王夫人又命人打听，都说是好。}
}
{
    \eR{[English source not present in the checked-in Bencraft/Joly files.]}
}
{
}

\Verse{187}
{
    \zR{王夫人便 告诉了宝钗，仍请了薛姨妈细细的告诉了袭人。}
}
{
    \eR{[English source not present in the checked-in Bencraft/Joly files.]}
}
{
}

\Verse{188}
{
    \zR{袭人悲伤不已，又不敢违命的，心 里想起宝玉那年到他家去，回来说的死也不回去的话，“如今太太硬作主张，若说
        我守着，又叫人说我不害臊；若是去了，实不是我的心愿。”}
}
{
    \eR{[English source not present in the checked-in Bencraft/Joly files.]}
}
{
}

\Verse{189}
{
    \zR{便哭得咽哽难鸣。}
}
{
    \eR{[English source not present in the checked-in Bencraft/Joly files.]}
}
{
}

\Verse{190}
{
    \zR{又 被薛姨妈宝钗等苦劝，回过念头想道：“我若是死在这里，倒把太太的好心弄坏了， 我该死在家里才是。”}
}
{
    \eR{[English source not present in the checked-in Bencraft/Joly files.]}
}
{
}

\Verse{191}
{
    \zR{于是袭人含悲叩辞了众人。}
}
{
    \eR{[English source not present in the checked-in Bencraft/Joly files.]}
}
{
}

\Verse{192}
{
    \zR{那姐妹分手时，自然更有一番不 忍说。}
}
{
    \eR{[English source not present in the checked-in Bencraft/Joly files.]}
}
{
}

\Verse{193}
{
    \zR{袭人怀着必死的心肠，上车回去，见了哥哥嫂子，也是哭泣，但只说不出来。}
}
{
    \eR{[English source not present in the checked-in Bencraft/Joly files.]}
}
{
}

\Verse{194}
{
    \zR{那花自芳悉把蒋家的聘礼送给他看，又把自己所办妆奁一一指给他瞧，说：“那是 太太赏的，那是置办的。”}
}
{
    \eR{[English source not present in the checked-in Bencraft/Joly files.]}
}
{
}

\Verse{195}
{
    \zR{袭人此时更难开口。}
}
{
    \eR{[English source not present in the checked-in Bencraft/Joly files.]}
}
{
}

\Verse{196}
{
    \zR{住了两天，细想起来：“哥哥办事 不错。}
}
{
    \eR{[English source not present in the checked-in Bencraft/Joly files.]}
}
{
}

\Verse{197}
{
    \zR{若是死在哥哥家里，岂不又害了哥哥呢？”}
}
{
    \eR{[English source not present in the checked-in Bencraft/Joly files.]}
}
{
}

\Verse{198}
{
    \zR{千思万想，左右为难，真是一缕 柔肠，几乎牵断，只得忍住。}
}
{
    \eR{[English source not present in the checked-in Bencraft/Joly files.]}
}
{
}

\Verse{199}
{
    \zR{那日已是迎娶吉期，袭人本不是那一种泼辣人，委委屈屈的上轿而去，心里另 想到那里再作打算。}
}
{
    \eR{[English source not present in the checked-in Bencraft/Joly files.]}
}
{
}

\Verse{200}
{
    \zR{岂知过了门，见那蒋家办事，极其认真，全都按着正配的规矩。}
}
{
    \eR{[English source not present in the checked-in Bencraft/Joly files.]}
}
{
}

\Verse{201}
{
    \zR{一进了门，丫头仆妇，都称“奶奶”。}
}
{
    \eR{[English source not present in the checked-in Bencraft/Joly files.]}
}
{
}

\Verse{202}
{
    \zR{袭人此时欲要死在这里，又恐害了人家，辜 负了一番好意。}
}
{
    \eR{[English source not present in the checked-in Bencraft/Joly files.]}
}
{
}

\Verse{203}
{
    \zR{那夜原是哭着不肯俯就的，那姑爷却极柔情曲意的承顺。}
}
{
    \eR{[English source not present in the checked-in Bencraft/Joly files.]}
}
{
}

\Verse{204}
{
    \zR{到了第二 天开箱，这姑爷看见一条猩红汗巾，方知是宝玉的丫头。}
}
{
    \eR{[English source not present in the checked-in Bencraft/Joly files.]}
}
{
}

\Verse{205}
{
    \zR{原来当初只知是贾母的侍 儿，益想不到是袭人。}
}
{
    \eR{[English source not present in the checked-in Bencraft/Joly files.]}
}
{
}

\Verse{206}
{
    \zR{此时蒋玉函念着宝玉待他的旧情，倒觉满心惶愧，更加周旋； 又故意将宝玉所换那条松花绿的汗巾拿出来。}
}
{
    \eR{[English source not present in the checked-in Bencraft/Joly files.]}
}
{
}

\Verse{207}
{
    \zR{袭人看了，方知这姓蒋的原来就是蒋 玉函，始信姻缘前定。}
}
{
    \eR{[English source not present in the checked-in Bencraft/Joly files.]}
}
{
}

\Verse{208}
{
    \zR{袭人才将心事说出。}
}
{
    \eR{[English source not present in the checked-in Bencraft/Joly files.]}
}
{
}

\Verse{209}
{
    \zR{蒋玉函也深为叹息敬服，不敢勉强，并 越发温柔体贴，弄得个袭人真无死所了。}
}
{
    \eR{[English source not present in the checked-in Bencraft/Joly files.]}
}
{
}

\Verse{210}
{
    \zR{看官听说：虽然事有前定，无可奈何，但 孽子孤臣，义夫节妇，这“不得已”三字也不是一概推委得的。}
}
{
    \eR{[English source not present in the checked-in Bencraft/Joly files.]}
}
{
}

\Verse{211}
{
    \zR{此袭人所以在“又 副册”也。}
}
{
    \eR{[English source not present in the checked-in Bencraft/Joly files.]}
}
{
}

\Verse{212}
{
    \zR{正是前人过那桃花庙的诗上说道： 千古艰难惟一死，伤心岂独息夫人! 不言袭人从此又是一番天地。}
}
{
    \eR{[English source not present in the checked-in Bencraft/Joly files.]}
}
{
}

\Verse{213}
{
    \zR{且说那贾雨村犯了婪索的案件，审明定罪，今遇 大赦，递籍为民。}
}
{
    \eR{[English source not present in the checked-in Bencraft/Joly files.]}
}
{
}

\Verse{214}
{
    \zR{雨村因叫家眷先行，自己带了一个小厮，一车行李，来到急流津 觉迷渡口。}
}
{
    \eR{[English source not present in the checked-in Bencraft/Joly files.]}
}
{
}

\Verse{215}
{
    \zR{只见一个道者，从那渡头草棚里出来，执手相迎。}
}
{
    \eR{[English source not present in the checked-in Bencraft/Joly files.]}
}
{
}

\Verse{216}
{
    \zR{雨村认得是甄士隐， 也连忙打恭。}
}
{
    \eR{[English source not present in the checked-in Bencraft/Joly files.]}
}
{
}

\Verse{217}
{
    \zR{士隐道：“贾老先生，别来无恙？”}
}
{
    \eR{[English source not present in the checked-in Bencraft/Joly files.]}
}
{
}

\Verse{218}
{
    \zR{雨村道：“老仙长到底是甄老先 生!何前次相逢，觌面不认？}
}
{
    \eR{[English source not present in the checked-in Bencraft/Joly files.]}
}
{
}

\Verse{219}
{
    \zR{后知火焚草亭，鄙下深为惶恐。}
}
{
    \eR{[English source not present in the checked-in Bencraft/Joly files.]}
}
{
}

\Verse{220}
{
    \zR{今日幸得相逢，益叹老 仙翁道德高深。}
}
{
    \eR{[English source not present in the checked-in Bencraft/Joly files.]}
}
{
}

\Verse{221}
{
    \zR{奈鄙人下愚不移，致有今日。”}
}
{
    \eR{[English source not present in the checked-in Bencraft/Joly files.]}
}
{
}

\Verse{222}
{
    \zR{甄士隐道：“前者老大人高官显爵， 贫道怎敢相认？}
}
{
    \eR{[English source not present in the checked-in Bencraft/Joly files.]}
}
{
}

\Verse{223}
{
    \zR{原因故交，敢赠片言，不意老大人相弃之深。}
}
{
    \eR{[English source not present in the checked-in Bencraft/Joly files.]}
}
{
}

\Verse{224}
{
    \zR{然而富贵穷通，亦非 偶然，今日复得相逢，也是一桩奇事。}
}
{
    \eR{[English source not present in the checked-in Bencraft/Joly files.]}
}
{
}

\Verse{225}
{
    \zR{这里离草庵不远，暂请膝谈，未知可否？”}
}
{
    \eR{[English source not present in the checked-in Bencraft/Joly files.]}
}
{
}

\Verse{226}
{
    \zR{雨村欣然领命。}
}
{
    \eR{[English source not present in the checked-in Bencraft/Joly files.]}
}
{
}

\Verse{227}
{
    \zR{两人携手而行，小厮驱车随后，到了一座茅庵。}
}
{
    \eR{[English source not present in the checked-in Bencraft/Joly files.]}
}
{
}

\Verse{228}
{
    \zR{士隐让进，雨村坐下，小童献 茶上来。}
}
{
    \eR{[English source not present in the checked-in Bencraft/Joly files.]}
}
{
}

\Verse{229}
{
    \zR{雨村便请教仙长超尘始末。}
}
{
    \eR{[English source not present in the checked-in Bencraft/Joly files.]}
}
{
}

\Verse{230}
{
    \zR{士隐笑道：“一念之间，尘凡顿易。}
}
{
    \eR{[English source not present in the checked-in Bencraft/Joly files.]}
}
{
}

\Verse{231}
{
    \zR{老先生从 繁华境中来，岂不知温柔富贵乡中有一宝玉乎？”}
}
{
    \eR{[English source not present in the checked-in Bencraft/Joly files.]}
}
{
}

\Verse{232}
{
    \zR{雨村道：“怎么不知。}
}
{
    \eR{[English source not present in the checked-in Bencraft/Joly files.]}
}
{
}

\Verse{233}
{
    \zR{近闻纷纷 传述，说他也遁入空门。}
}
{
    \eR{[English source not present in the checked-in Bencraft/Joly files.]}
}
{
}

\Verse{234}
{
    \zR{下愚当时也曾与他往来过数次，再不想此人竟有如是之决 绝。”}
}
{
    \eR{[English source not present in the checked-in Bencraft/Joly files.]}
}
{
}

\Verse{235}
{
    \zR{士隐道：“非也。}
}
{
    \eR{[English source not present in the checked-in Bencraft/Joly files.]}
}
{
}

\Verse{236}
{
    \zR{这一段奇缘，我先知之。}
}
{
    \eR{[English source not present in the checked-in Bencraft/Joly files.]}
}
{
}

\Verse{237}
{
    \zR{昔年我与先生在仁清巷旧宅门口 叙话之前，我已会过他一面。”}
}
{
    \eR{[English source not present in the checked-in Bencraft/Joly files.]}
}
{
}

\Verse{238}
{
    \zR{雨村惊讶道：“京城离贵乡甚远，何以能见？”}
}
{
    \eR{[English source not present in the checked-in Bencraft/Joly files.]}
}
{
}

\Verse{239}
{
    \zR{士 隐道：“神交久矣。”}
}
{
    \eR{[English source not present in the checked-in Bencraft/Joly files.]}
}
{
}

\Verse{240}
{
    \zR{雨村道：“既然如此，现今宝玉的下落，仙长定能知之？”}
}
{
    \eR{[English source not present in the checked-in Bencraft/Joly files.]}
}
{
}

\Verse{241}
{
    \zR{士隐道：“宝玉，即‘宝玉’也。}
}
{
    \eR{[English source not present in the checked-in Bencraft/Joly files.]}
}
{
}

\Verse{242}
{
    \zR{那年荣宁查抄之前，钗黛分离之日，此玉早已离 世：一为避祸，二为撮合。}
}
{
    \eR{[English source not present in the checked-in Bencraft/Joly files.]}
}
{
}

\Verse{243}
{
    \zR{从此夙缘一了，形质归一。}
}
{
    \eR{[English source not present in the checked-in Bencraft/Joly files.]}
}
{
}

\Verse{244}
{
    \zR{又复稍示神灵，高魁贵子， 方显得此玉乃天奇地灵锻炼之宝，非凡间可比。}
}
{
    \eR{[English source not present in the checked-in Bencraft/Joly files.]}
}
{
}

\Verse{245}
{
    \zR{前经茫茫大士渺渺真人携带下凡， 如今尘缘已满，仍是此二人携归本处：便是宝玉的下落。”}
}
{
    \eR{[English source not present in the checked-in Bencraft/Joly files.]}
}
{
}

\Verse{246}
{
    \zR{雨村听了，虽不能全然 明白，却也十知四五，便点头叹道：“原来如此，下愚不知。}
}
{
    \eR{[English source not present in the checked-in Bencraft/Joly files.]}
}
{
}

\Verse{247}
{
    \zR{但那宝玉既有如此的 来历，又何以情迷至此，复又豁悟如此？}
}
{
    \eR{[English source not present in the checked-in Bencraft/Joly files.]}
}
{
}

\Verse{248}
{
    \zR{还要请教。”}
}
{
    \eR{[English source not present in the checked-in Bencraft/Joly files.]}
}
{
}

\Verse{249}
{
    \zR{士隐笑道：“此事说来，先 生未必尽解。}
}
{
    \eR{[English source not present in the checked-in Bencraft/Joly files.]}
}
{
}

\Verse{250}
{
    \zR{太虚幻境，即是真如福地。}
}
{
    \eR{[English source not present in the checked-in Bencraft/Joly files.]}
}
{
}

\Verse{251}
{
    \zR{两番阅册，原始要终之道，历历生平，如 何不悟？}
}
{
    \eR{[English source not present in the checked-in Bencraft/Joly files.]}
}
{
}

\Verse{252}
{
    \zR{仙草归真，焉有通灵不复原之理呢？”}
}
{
    \eR{[English source not present in the checked-in Bencraft/Joly files.]}
}
{
}

\Verse{253}
{
    \zR{雨村听着，却不明白，知是仙机，也不便更问。}
}
{
    \eR{[English source not present in the checked-in Bencraft/Joly files.]}
}
{
}

\Verse{254}
{
    \zR{因又说道：“宝玉之事，既得 闻命。}
}
{
    \eR{[English source not present in the checked-in Bencraft/Joly files.]}
}
{
}

\Verse{255}
{
    \zR{但敝族闺秀如是之多，何元妃以下，算来结局俱属平常呢？”}
}
{
    \eR{[English source not present in the checked-in Bencraft/Joly files.]}
}
{
}

\Verse{256}
{
    \zR{士隐叹道：“老 先生莫怪拙言!贵族之女，俱属从情天孽海而来。}
}
{
    \eR{[English source not present in the checked-in Bencraft/Joly files.]}
}
{
}

\Verse{257}
{
    \zR{大凡古今女子，那‘淫’字固不 可犯，只这‘情’字也是沾染不得的。}
}
{
    \eR{[English source not present in the checked-in Bencraft/Joly files.]}
}
{
}

\Verse{258}
{
    \zR{所以崔莺苏小，无非仙子尘心；宋玉相如， 大是文人口孽。}
}
{
    \eR{[English source not present in the checked-in Bencraft/Joly files.]}
}
{
}

\Verse{259}
{
    \zR{但凡情思缠绵，那结局就不可问了。”}
}
{
    \eR{[English source not present in the checked-in Bencraft/Joly files.]}
}
{
}

\Verse{260}
{
    \zR{雨村听到这里，不觉拈须长叹。}
}
{
    \eR{[English source not present in the checked-in Bencraft/Joly files.]}
}
{
}

\Verse{261}
{
    \zR{因又问道：“请教仙翁：那荣宁两府，尚可如 前否？”}
}
{
    \eR{[English source not present in the checked-in Bencraft/Joly files.]}
}
{
}

\Verse{262}
{
    \zR{士隐道：“福善祸淫，古今定理。}
}
{
    \eR{[English source not present in the checked-in Bencraft/Joly files.]}
}
{
}

\Verse{263}
{
    \zR{现今荣宁两府，善者修缘，恶者悔祸， 将来兰桂齐芳，家道复初，也是自然的道理。”}
}
{
    \eR{[English source not present in the checked-in Bencraft/Joly files.]}
}
{
}

\Verse{264}
{
    \zR{雨村低了半日头，忽然笑道：“是 了，是了。}
}
{
    \eR{[English source not present in the checked-in Bencraft/Joly files.]}
}
{
}

\Verse{265}
{
    \zR{现在他府中有一个名兰的，已中乡榜，恰好应着‘兰’字。}
}
{
    \eR{[English source not present in the checked-in Bencraft/Joly files.]}
}
{
}

\Verse{266}
{
    \zR{适间老仙翁 说‘兰桂齐芳’，又道‘宝玉高魁贵子’，莫非他有遗腹之子，可以飞黄腾达的么？”}
}
{
    \eR{[English source not present in the checked-in Bencraft/Joly files.]}
}
{
}

\Verse{267}
{
    \zR{士隐微微笑道：“此系后事，未便预说。”}
}
{
    \eR{[English source not present in the checked-in Bencraft/Joly files.]}
}
{
}

\Verse{268}
{
    \zR{雨村还要再问，士隐不答，便命人设具盘飧，邀雨村共食。}
}
{
    \eR{[English source not present in the checked-in Bencraft/Joly files.]}
}
{
}

\Verse{269}
{
    \zR{食毕，雨村还要问 自己的终身。}
}
{
    \eR{[English source not present in the checked-in Bencraft/Joly files.]}
}
{
}

\Verse{270}
{
    \zR{士隐便道：“老先生草庵暂歇。}
}
{
    \eR{[English source not present in the checked-in Bencraft/Joly files.]}
}
{
}

\Verse{271}
{
    \zR{我还有一段俗缘未了，正当今日完结。”}
}
{
    \eR{[English source not present in the checked-in Bencraft/Joly files.]}
}
{
}

\Verse{272}
{
    \zR{雨村惊讶道：“仙长纯修若此，不知尚有何俗缘？”}
}
{
    \eR{[English source not present in the checked-in Bencraft/Joly files.]}
}
{
}

\Verse{273}
{
    \zR{士隐道：“也不过是儿女私情 罢了。”}
}
{
    \eR{[English source not present in the checked-in Bencraft/Joly files.]}
}
{
}

\Verse{274}
{
    \zR{雨村听了，益发惊异：“请问仙长何出此言？”}
}
{
    \eR{[English source not present in the checked-in Bencraft/Joly files.]}
}
{
}

\Verse{275}
{
    \zR{士隐道：“老先生有所不 知：小女英莲，幼遭尘劫，老先生初任之时，曾经判断。}
}
{
    \eR{[English source not present in the checked-in Bencraft/Joly files.]}
}
{
}

\Verse{276}
{
    \zR{今归薛姓，产难完劫，遗 一子于薛家，以承宗祧。}
}
{
    \eR{[English source not present in the checked-in Bencraft/Joly files.]}
}
{
}

\Verse{277}
{
    \zR{此时正是尘缘脱尽之时，只好接引接引。”}
}
{
    \eR{[English source not present in the checked-in Bencraft/Joly files.]}
}
{
}

\Verse{278}
{
    \zR{士隐说着，拂 袖而起。}
}
{
    \eR{[English source not present in the checked-in Bencraft/Joly files.]}
}
{
}

\Verse{279}
{
    \zR{雨村心中恍恍惚惚，就在这急流津觉迷渡口草庵中睡着了。}
}
{
    \eR{[English source not present in the checked-in Bencraft/Joly files.]}
}
{
}

\Verse{280}
{
    \zR{这士隐自去度脱了香菱，送到太虚幻境，交那警幻仙子对册。}
}
{
    \eR{[English source not present in the checked-in Bencraft/Joly files.]}
}
{
}

\Verse{281}
{
    \zR{刚过牌坊，见那 一僧一道缥缈而来，士隐接着说道：“大士、真人，恭喜贺喜!情缘完结，都交割 清楚了么？”}
}
{
    \eR{[English source not present in the checked-in Bencraft/Joly files.]}
}
{
}

\Verse{282}
{
    \zR{那僧道说：“情缘尚未全结，倒是那蠢物已经回来了。}
}
{
    \eR{[English source not present in the checked-in Bencraft/Joly files.]}
}
{
}

\Verse{283}
{
    \zR{还得把他送还 原所，将他的后事叙明，不枉他下世一回。”}
}
{
    \eR{[English source not present in the checked-in Bencraft/Joly files.]}
}
{
}

\Verse{284}
{
    \zR{士隐听了，便拱手而别。}
}
{
    \eR{[English source not present in the checked-in Bencraft/Joly files.]}
}
{
}

\Verse{285}
{
    \zR{那僧道仍携 了玉到青埂峰下，将“宝玉”安放在女娲炼石补天之处，各自云游而去。}
}
{
    \eR{[English source not present in the checked-in Bencraft/Joly files.]}
}
{
}

\Verse{286}
{
    \zR{从此后： 天外书传天外事，两番人作一番人。}
}
{
    \eR{[English source not present in the checked-in Bencraft/Joly files.]}
}
{
}

\Verse{287}
{
    \zR{这一日，空空道人又从青埂峰前经过，见那补天未用之石仍在那里，上面字迹 依然如旧，又从头的细细看了一遍。}
}
{
    \eR{[English source not present in the checked-in Bencraft/Joly files.]}
}
{
}

\Verse{288}
{
    \zR{见后面偈文后又历叙了多少收缘结果的话头， 便点头叹道：“我从前见石兄这段奇文，原说可以闻世传奇，所以曾经抄录，但未 见返本还原。}
}
{
    \eR{[English source not present in the checked-in Bencraft/Joly files.]}
}
{
}

\Verse{289}
{
    \zR{不知何时，复有此段佳话？}
}
{
    \eR{[English source not present in the checked-in Bencraft/Joly files.]}
}
{
}

\Verse{290}
{
    \zR{方知石兄下凡一次，磨出光明，修成圆觉， 也可谓无复遗憾了。}
}
{
    \eR{[English source not present in the checked-in Bencraft/Joly files.]}
}
{
}

\Verse{291}
{
    \zR{只怕年深日久，字迹模糊，反有舛错，不如我再抄录一番，寻 个世上清闲无事的人，托他传遍，知道奇而不奇，俗而不俗，真而不真，假而不假。}
}
{
    \eR{[English source not present in the checked-in Bencraft/Joly files.]}
}
{
}

\Verse{292}
{
    \zR{或者尘梦劳人，聊倩鸟呼归去；山灵好客，更从石化飞来：亦未可知。”}
}
{
    \eR{[English source not present in the checked-in Bencraft/Joly files.]}
}
{
}

\Verse{293}
{
    \zR{想毕，便 又抄了，仍袖至那繁华昌盛地方。}
}
{
    \eR{[English source not present in the checked-in Bencraft/Joly files.]}
}
{
}

\Verse{294}
{
    \zR{遍寻了一番，不是建功立业之人，即系糊口谋衣 之辈，那有闲情去和石头饶舌？}
}
{
    \eR{[English source not present in the checked-in Bencraft/Joly files.]}
}
{
}

\Verse{295}
{
    \zR{直寻到急流津觉迷渡口草庵中，睡着一个人，因想 他必是闲人，便要将这抄录的《石头记》给他看看。}
}
{
    \eR{[English source not present in the checked-in Bencraft/Joly files.]}
}
{
}

\Verse{296}
{
    \zR{那知那人再叫不醒。}
}
{
    \eR{[English source not present in the checked-in Bencraft/Joly files.]}
}
{
}

\Verse{297}
{
    \zR{空空道人 复又使劲拉他，才慢慢的开眼坐起。}
}
{
    \eR{[English source not present in the checked-in Bencraft/Joly files.]}
}
{
}

\Verse{298}
{
    \zR{便接来草草一看，仍旧掷下道：“这事我已亲 见尽知，你这抄录的尚无舛错。}
}
{
    \eR{[English source not present in the checked-in Bencraft/Joly files.]}
}
{
}

\Verse{299}
{
    \zR{我只指与你一个人，托他传去，便可归结这段新鲜 公案了。”}
}
{
    \eR{[English source not present in the checked-in Bencraft/Joly files.]}
}
{
}

\Verse{300}
{
    \zR{空空道人忙问何人，那人道：“你须待某年某月某日某时，到一个悼红 轩中，有个曹雪芹先生。}
}
{
    \eR{[English source not present in the checked-in Bencraft/Joly files.]}
}
{
}

\Verse{301}
{
    \zR{只说贾雨村言，托他如此如此。”}
}
{
    \eR{[English source not present in the checked-in Bencraft/Joly files.]}
}
{
}

\Verse{302}
{
    \zR{说毕，仍旧睡下了。}
}
{
    \eR{[English source not present in the checked-in Bencraft/Joly files.]}
}
{
}

\Verse{303}
{
    \zR{那空空道人牢牢记着此言，又不知过了几世几劫，果然有个悼红轩，见那曹雪 芹先生正在那里翻阅历来的古史。}
}
{
    \eR{[English source not present in the checked-in Bencraft/Joly files.]}
}
{
}

\Verse{304}
{
    \zR{空空道人便将贾雨村言了，方把这《石头记》示 看。}
}
{
    \eR{[English source not present in the checked-in Bencraft/Joly files.]}
}
{
}

\Verse{305}
{
    \zR{那雪芹先生笑道：“果然是‘贾雨村言’了！”}
}
{
    \eR{[English source not present in the checked-in Bencraft/Joly files.]}
}
{
}

\Verse{306}
{
    \zR{空空道人便问：“先生何以认 得此人，便肯替他传述？”}
}
{
    \eR{[English source not present in the checked-in Bencraft/Joly files.]}
}
{
}

\Verse{307}
{
    \zR{那雪芹先生笑道：“说你‘空空’，原来肚里果然空空。}
}
{
    \eR{[English source not present in the checked-in Bencraft/Joly files.]}
}
{
}

\Verse{308}
{
    \zR{既是‘假语村言’，但无鲁鱼亥豕以及背谬矛盾之处，乐得与二三同志，酒馀饭饱， 雨夕灯窗，同消寂寞，又不必大人先生品题传世。}
}
{
    \eR{[English source not present in the checked-in Bencraft/Joly files.]}
}
{
}

\Verse{309}
{
    \zR{似你这样寻根究底，便是刻舟求 剑、胶柱鼓瑟了。”}
}
{
    \eR{[English source not present in the checked-in Bencraft/Joly files.]}
}
{
}

\Verse{310}
{
    \zR{那空空道人听了，仰天大笑，掷下抄本，飘然而去。}
}
{
    \eR{[English source not present in the checked-in Bencraft/Joly files.]}
}
{
}

\Verse{311}
{
    \zR{一面走着， 口中说道：“原来是敷衍荒唐!不但作者不知，抄者不知，并阅者也不知。}
}
{
    \eR{[English source not present in the checked-in Bencraft/Joly files.]}
}
{
}

\Verse{312}
{
    \zR{不过游 戏笔墨，陶情适性而已！”}
}
{
    \eR{[English source not present in the checked-in Bencraft/Joly files.]}
}
{
}

\Verse{313}
{
    \zR{后人见了这本传奇，亦曾题过四句偈语，为作者缘起之言更进一竿。}
}
{
    \eR{[English source not present in the checked-in Bencraft/Joly files.]}
}
{
}

\Verse{314}
{
    \zR{云： 说到辛酸处，荒唐愈可悲。}
}
{
    \eR{[English source not present in the checked-in Bencraft/Joly files.]}
}
{
}

\Verse{315}
{
    \zR{由来同一梦，休笑世人痴! (本章完)}
}
{
    \eR{[English source not present in the checked-in Bencraft/Joly files.]}
}
{
}
